% ============================================================================
% Chapter 12: Software - Python
% Book source: part1-linear-programming/ch03-software/software-python-book1.tex
% Exercise numbering synced to \begin{ex} order as of 2026-07-13.
% All code executed with Python 3.10, PuLP 3.3.2, bundled CBC solver.
% ============================================================================

\section{Bab 12: Perangkat Lunak -- Python}

Bab ini memiliki sembilan latihan yang semuanya dibangun berdasarkan dua
contoh PuLP berkelanjutan pada bagian PuLP dalam bab (Bagian 12.12): LP bauran
produk dan LP transportasi lengkap. Latihan pemanasan (12.1--12.3) mengubah
data pada contoh yang telah dikerjakan; soal inti (12.4--12.6) membangun model
dari awal, mengekstrak dual dan slack, serta menjalankan sapuan kendala
$\varepsilon$; soal konsep (12.7--12.8) mendiagnosis dua kutu pemodelan PuLP
yang klasik (kendala yang hilang dan arah optimisasi yang salah); sedangkan
soal tantangan (12.9) menelusuri laba optimal sebagai fungsi ruas kanan.
Setiap solusi di bawah ini dihasilkan dengan benar-benar menjalankan kode yang
ditampilkan (PuLP 3.3.2 dengan pemecah CBC bawaan); keluaran tercetak disalin
persis. Solusi Terpilih buku untuk 12.2, 12.5, dan 12.9 dijalankan ulang dan
dikonfirmasi; solusi tersebut digunakan kembali dan diperluas di sini.

% ----------------------------------------------------------------------------
\exsol{12.1}{Margin Laba Baru}{1}

Hanya baris fungsi tujuan yang berubah: koefisien laba $X_1$ menjadi 5.

\begin{verbatim}
from pulp import *

prob = LpProblem("Bauran_Produk", LpMaximize)
x1 = LpVariable("X1", lowBound=0)
x2 = LpVariable("X2", lowBound=0)

prob += 5*x1 + 2*x2, "Laba_Total"       # sebelumnya 3*x1 + 2*x2
prob += 10*x1 + 5*x2 <= 300, "Bahan"
prob += 4*x1 + 4*x2 <= 160, "Tenaga_Kerja"
prob += 2*x1 + 6*x2 <= 180, "Mesin"

prob.solve(PULP_CBC_CMD(msg=0))
print(f"Status: {LpStatus[prob.status]}")
print(f"Laba Maksimum: ${value(prob.objective)}")
print(f"X1 = {x1.varValue}")
print(f"X2 = {x2.varValue}")
for name, c in prob.constraints.items():
    print(name, "kelonggaran =", c.slack)
\end{verbatim}

Keluaran sebenarnya:

\begin{verbatim}
Status: Optimal
Laba Maksimum: $150.0
X1 = 30.0
X2 = 0.0
Bahan kelonggaran = -0.0
Tenaga_Kerja kelonggaran = 40.0
Mesin kelonggaran = 120.0
\end{verbatim}

Rencana baru adalah $(X_1, X_2) = (30, 0)$ dengan laba $\$150$. Produk 1 kini
memberikan laba per unit bahan langka yang begitu tinggi (laba $5$ per $10$
unit bahan, dibandingkan $2$ per $5$ unit untuk Produk 2, dan bahan merupakan
sumber daya yang pertama kali habis) sehingga perusahaan sepenuhnya
berspesialisasi. Hanya kendala \textbf{Bahan} yang mengikat
($10 \cdot 30 + 5 \cdot 0 = 300$); Tenaga Kerja memiliki slack $40$ dan Mesin
memiliki slack $120$ (batas nonnegativitas $X_2 \geq 0$ juga aktif). Sesuai
yang dijanjikan, jawabannya bukan lagi $(20,20)$.

% ----------------------------------------------------------------------------
\exsol{12.2}{Dari Matematika ke PuLP}{1}

Solusi Terpilih dalam buku sudah benar sebagaimana tercetak; kode di bawah
dijalankan ulang untuk mengonfirmasinya. Dengan mengikuti pola enam langkah:

\begin{verbatim}
from pulp import *

# Langkah 1: masalah
prob = LpProblem("Three_Var", LpMaximize)
# Langkah 2: variabel
x1 = LpVariable("X1", lowBound=0)
x2 = LpVariable("X2", lowBound=0)
x3 = LpVariable("X3", lowBound=0)

# Langkah 3: objektif
prob += 4*x1 + 6*x2 + 5*x3, "Objective"
prob += 2*x1 + x2 + x3 <= 9, "C1"                # Langkah 4: kendala
prob += x1 + 3*x2 + 2*x3 <= 14, "C2"
prob += x1 + x2 + 2*x3 <= 10, "C3"

# Langkah 5: selesaikan
prob.solve(PULP_CBC_CMD(msg=0))
print(LpStatus[prob.status], value(prob.objective))
for v in prob.variables():                         # Langkah 6: hasil
    print(v.name, "=", v.varValue)
\end{verbatim}

Keluaran sebenarnya:

\begin{verbatim}
Optimal 35.0
X1 = 2.0
X2 = 2.0
X3 = 3.0
\end{verbatim}

Nilai tujuan optimal adalah $35$ pada $(x_1, x_2, x_3) = (2, 2, 3)$.
Ketiga kendala mengikat ($4+2+3 = 9$, $2+6+6 = 14$, $2+2+6 = 10$);

\begin{qaflag}
Diselesaikan 2026-07-10: data semula untuk latihan ini cocok koefisien demi
koefisien dengan sebuah contoh buku ajar berhak cipta yang terkenal. Data
tersebut diganti dengan kasus orisinal (optimum terverifikasi $(2,2,3)$,
nilai $35$), baik dalam buku maupun dalam manual ini.
\end{qaflag}

% ----------------------------------------------------------------------------
\exsol{12.3}{Data Permintaan Baru}{1}

Hanya kamus \texttt{demand} yang berubah.

\begin{verbatim}
from pulp import *

plants = ['Chicago', 'Denver']
markets = ['NYC', 'LA', 'Houston']

supply = {'Chicago': 100, 'Denver': 150}
# Permintaan diperbarui
demand = {'NYC': 100, 'LA': 80, 'Houston': 70}

cost = {
    ('Chicago', 'NYC'): 3,
    ('Chicago', 'LA'): 5,
    ('Chicago', 'Houston'): 2,
    ('Denver', 'NYC'): 4,
    ('Denver', 'LA'): 2,
    ('Denver', 'Houston'): 3,
}

prob = LpProblem("Transportation", LpMinimize)
routes = [(p, m) for p in plants for m in markets]
x = LpVariable.dicts("Ship", routes, lowBound=0)
prob += lpSum([cost[r] * x[r] for r in routes]), "Total_Cost"
for p in plants:
    prob += (
        lpSum([x[p,m] for m in markets]) <= supply[p],
        f"Supply_{p}",
    )
for m in markets:
    prob += (
        lpSum([x[p,m] for p in plants]) >= demand[m],
        f"Demand_{m}",
    )

prob.solve(PULP_CBC_CMD(msg=0))
print(f"Status: {LpStatus[prob.status]}")
print(f"Biaya Total: ${value(prob.objective)}")
for p in plants:
    for m in markets:
        if x[p,m].varValue > 0:
            print(f"  {p} -> {m}: {x[p,m].varValue}")
\end{verbatim}

Keluaran sebenarnya:

\begin{verbatim}
Status: Optimal
Biaya Total: $670.0
  Chicago -> NYC: 100.0
  Denver -> LA: 80.0
  Denver -> Houston: 70.0
\end{verbatim}

Biaya total baru adalah $\$670$ (naik dari $\$610$). Rencana CBC mengirimkan
seluruh persediaan Chicago ke NYC dan memenuhi LA serta Houston dari Denver.
(Ada optimum alternatif dengan biaya yang sama, misalnya Denver melayani LA
dan NYC sementara Chicago memenuhi Houston serta sebagian NYC; mahasiswa
dapat melaporkan rencana berbeda dengan biaya $670$.)

Penjelasan satu kalimat: permintaan total meningkat dari $240$ menjadi $250$
unit dan bergeser dari LA (yang dilayani Denver dengan tarif termurah dalam
sistem, $\$2$) menuju NYC (tarif masuk termurah $\$3$), sehingga jaringan
harus mengirim lebih banyak unit melalui rute yang lebih mahal. Perhatikan
pula bahwa permintaan kini tepat sama dengan persediaan total ($250 = 250$),
sehingga tidak ada kapasitas cadangan di kedua pabrik yang dapat dimanfaatkan
untuk memilih rute yang lebih murah.

% ----------------------------------------------------------------------------
\exsol{12.4}{Model Transportasi Tiga Pabrik}{2}

\begin{verbatim}
from pulp import *

# === DATA ===
plants = ['A', 'B', 'C']
markets = ['M1', 'M2', 'M3']
supply = {'A': 60, 'B': 70, 'C': 80}
demand = {'M1': 50, 'M2': 60, 'M3': 70}
cost = {
    ('A','M1'): 4, ('A','M2'): 6, ('A','M3'): 8,
    ('B','M1'): 5, ('B','M2'): 4, ('B','M3'): 3,
    ('C','M1'): 6, ('C','M2'): 7, ('C','M3'): 5
}

# === MODEL ===
prob = LpProblem("Three_Plant_Transport", LpMinimize)
routes = [(p, m) for p in plants for m in markets]
x = LpVariable.dicts("Ship", routes, lowBound=0)

prob += lpSum([cost[r] * x[r] for r in routes]), "Biaya_Total"
for p in plants:
    prob += (
        lpSum([x[p,m] for m in markets]) <= supply[p],
        f"Supply_{p}",
    )
for m in markets:
    prob += (
        lpSum([x[p,m] for p in plants]) >= demand[m],
        f"Demand_{m}",
    )

# === PENYELESAIAN / HASIL ===
prob.solve(PULP_CBC_CMD(msg=0))
print(f"Status: {LpStatus[prob.status]}")
print(f"Biaya Total: {value(prob.objective)}")
for p in plants:
    shipped = sum(x[p,m].varValue for m in markets)
    print(f"Pabrik {p}: mengirim {shipped} dari {supply[p]}"
          f" (tak terpakai {supply[p]-shipped})")
for p in plants:
    for m in markets:
        if x[p,m].varValue > 0:
            print(f"  {p} -> {m}: {x[p,m].varValue}")
\end{verbatim}

Keluaran sebenarnya:

\begin{verbatim}
Status: Optimal
Biaya Total: 770.0
Pabrik A: mengirim 50.0 dari 60 (tak terpakai 10.0)
Pabrik B: mengirim 70.0 dari 70 (tak terpakai 0.0)
Pabrik C: mengirim 60.0 dari 80 (tak terpakai 20.0)
  A -> M1: 50.0
  B -> M2: 60.0
  B -> M3: 10.0
  C -> M3: 60.0
\end{verbatim}

Biaya total minimum: $\$770$. Rencana: A memasok seluruh M1 (50 unit dengan
biaya 4), B memasok seluruh M2 (60 dengan biaya 4) ditambah 10 unit M3 (dengan
biaya 3), dan C memasok 60 unit M3 yang tersisa (dengan biaya 5). Pabrik A
dan C memiliki kapasitas tak terpakai (masing-masing 10 dan 20 unit); B, yang
rutenya menuju M2 dan M3 merupakan rute termurah pada kolom masing-masing,
digunakan sampai kapasitas penuh. Pemeriksaan kewajaran secara manual:
penetapan serakah M1$\leftarrow$A, M2$\leftarrow$B, M3$\leftarrow$B lalu C
menghasilkan $200 + 240 + 30 + 300 = 770$, dan LP memastikan tidak ada rencana
yang lebih murah.

% ----------------------------------------------------------------------------
\exsol{12.5}{Dual dan Slack dari Model yang Telah Diselesaikan}{2}

Skrip lengkap (data bauran produk semula, laba $3X_1 + 2X_2$):

\begin{verbatim}
from pulp import *

prob = LpProblem("Product_Mix", LpMaximize)
x1 = LpVariable("X1", lowBound=0)
x2 = LpVariable("X2", lowBound=0)
prob += 3*x1 + 2*x2, "Total_Profit"
prob += 10*x1 + 5*x2 <= 300, "Material"
prob += 4*x1 + 4*x2 <= 160, "Labor"
prob += 2*x1 + 6*x2 <= 180, "Machine"
prob.solve(PULP_CBC_CMD(msg=0))

print(
    LpStatus[prob.status], value(prob.objective),
    x1.varValue, x2.varValue,
)
for name, c in prob.constraints.items():
    print(name, "dual =", c.pi, "slack =", c.slack)
\end{verbatim}

Keluaran sebenarnya:

\begin{verbatim}
Optimal 100.0 20.0 20.0
Material dual = 0.2 slack = -0.0
Labor dual = 0.25 slack = -0.0
Machine dual = -0.0 slack = 20.0
\end{verbatim}

\begin{enumerate}
\item Bahan: dual $0.2$, slack $0$. Tenaga Kerja: dual $0.25$, slack $0$.
Mesin: dual $0$, slack $20$. (CBC mencetak $-0.0$; ini adalah nol titik
mengambang bertanda, yakni $0$.) Pemeriksaan dual secara manual: baris yang
mengikat memberikan $10y_1 + 4y_2 = 3$ dan $5y_1 + 4y_2 = 2$, sehingga
$y_1 = 0.2$, $y_2 = 0.25$.
\item Kendala Mesin memiliki nilai dual $0$ dan slack $20$: pada optimum
$(20,20)$, penggunaan mesin adalah $2(20) + 6(20) = 160 < 180$. Sumber daya
dengan kapasitas tersisa tidak bernilai pada margin, sehingga harga bayangannya
nol. Inilah kondisi kelonggaran komplementer: untuk setiap kendala, nilai dual atau slack
(atau keduanya) harus nol.
\item Nilai dual Tenaga Kerja adalah $0.25$, sehingga satu jam tenaga kerja
tambahan meningkatkan laba optimal sebesar $\$0.25$, dari $\$100$ menjadi
$\$100.25$ (berlaku selama basis saat ini tetap optimal; Latihan 12.9
menunjukkan bahwa laju ini berlaku hingga $b = 168$).
\end{enumerate}

% ----------------------------------------------------------------------------
\exsol{12.6}{Sapuan Kendala Epsilon}{2}

Dengan mengikuti templat dari bab multiobjektif (Bagian 13.3), kita
menyelesaikan kembali LP bauran produk dengan batas tambahan
$2X_1 + 6X_2 \leq \varepsilon$.

\begin{verbatim}
from pulp import *

print(f"{'eps':>4} {'X1':>6} {'X2':>6} {'laba':>7}")
for eps in [60, 90, 120, 150, 180]:
    prob = LpProblem("Product_Mix_Eps", LpMaximize)
    x1 = LpVariable("X1", lowBound=0)
    x2 = LpVariable("X2", lowBound=0)
    prob += 3*x1 + 2*x2
    prob += 10*x1 + 5*x2 <= 300, "Material"
    prob += 4*x1 + 4*x2 <= 160, "Labor"
    prob += 2*x1 + 6*x2 <= 180, "Machine"
    prob += 2*x1 + 6*x2 <= eps, "Machine_cap"
    prob.solve(PULP_CBC_CMD(msg=0))
    print(f"{eps:>4} {x1.varValue:>6.1f} {x2.varValue:>6.1f}"
          f" {value(prob.objective):>7.1f}")
\end{verbatim}

Keluaran sebenarnya:

\begin{verbatim}
 eps     X1     X2    laba
  60   30.0    0.0    90.0
  90   27.0    6.0    93.0
 120   24.0   12.0    96.0
 150   21.0   18.0    99.0
 180   20.0   20.0   100.0
\end{verbatim}

Komprominya: memperketat batas mesin mengurangi laba, tetapi biayanya kecil.
Dari $\varepsilon = 160$ (tempat optimum tanpa kendala tambahan $(20,20)$
menggunakan tepat $160$ jam mesin) hingga $\varepsilon = 60$, batas tersebut
mengikat dan setiap jam mesin yang dikorbankan mengurangi laba tepat
$\$0.10$ (harga bayangan batas): laba turun secara linear
$100 \to 99 \to 96 \to 93 \to 90$ ketika $\varepsilon$ turun
$160 \to 150 \to 120 \to 90 \to 60$. Untuk $\varepsilon \geq 160$
(termasuk $\varepsilon = 180$ yang tercetak), batas memiliki slack dan
solusinya adalah optimum semula. Sepanjang sapuan, rencana beralih secara
bertahap dari bauran yang sarat penggunaan mesin menuju Produk 1 saja:
perusahaan dapat memangkas penggunaan mesin hampir dua pertiga (dari $160$
menjadi $60$ jam) dengan kehilangan laba hanya $10\%$. Setiap baris merupakan
satu titik pada garis depan Pareto (laba, penggunaan mesin).

% ----------------------------------------------------------------------------
\exsol{12.7}{Kendala yang Tidak Ada}{2}

\begin{enumerate}
\item Perkiraan dan percobaan memberikan hasil yang sama: PuLP tidak
menimbulkan galat ataupun mencetak peringatan; PuLP dengan senang hati
menyelesaikan model yang Anda berikan. PuLP tidak dapat mengetahui bahwa suatu
kendala hilang; LP dengan dua kendala merupakan masalah yang sepenuhnya sah.
Eksekusi tersebut mengembalikan status \texttt{Optimal}.

\begin{verbatim}
from pulp import *

prob = LpProblem("Product_Mix_NoLabor", LpMaximize)
x1 = LpVariable("X1", lowBound=0)
x2 = LpVariable("X2", lowBound=0)
prob += 3*x1 + 2*x2, "Total_Profit"
prob += 10*x1 + 5*x2 <= 300, "Material"
# Baris tenaga kerja hilang
prob += 2*x1 + 6*x2 <= 180, "Machine"
prob.solve(PULP_CBC_CMD(msg=0))
print(
    LpStatus[prob.status], value(prob.objective),
    x1.varValue, x2.varValue,
)
print("Tenaga kerja terpakai:", 4*x1.varValue + 4*x2.varValue)
\end{verbatim}

Keluaran sebenarnya:
\begin{verbatim}
Optimal 102.0 18.0 24.0
Tenaga kerja terpakai: 168.0
\end{verbatim}

\item Solusi yang dikembalikan adalah $(18, 24)$ dengan laba $\$102$, yaitu
perpotongan garis Bahan dan Mesin ($10x_1 + 5x_2 = 300$,
$2x_1 + 6x_2 = 180$). Labanya melebihi $\$100$ karena menghapus suatu kendala
memperbesar daerah layak, sementara pada optimum sebenarnya $(20,20)$ kendala
Tenaga Kerja mengikat sehingga melonggarkannya benar-benar membantu. Rencana
ini \emph{tidak} dapat dilaksanakan: diperlukan
$4(18) + 4(24) = 168$ jam tenaga kerja, tetapi yang tersedia hanya $160$.

\item Dengan \texttt{prob += 4*x1 + 4*x2} (tanpa \texttt{<= 160}), PuLP
memperlakukan ekspresi itu sebagai \emph{fungsi tujuan} baru dan mengganti
$3x_1 + 2x_2$ dengan $4x_1 + 4x_2$. PuLP versi terbaru mencetak
\begin{verbatim}
UserWarning: Overwriting previously set objective.
\end{verbatim}
tetapi model tetap dapat diselesaikan dan melaporkan \texttt{Optimal}
(eksekusi kami menghasilkan $(18, 24)$ dengan nilai tujuan $168$, yaitu nilai
$4x_1 + 4x_2$, bukan laba). Kutu ini lebih berbahaya daripada baris yang
hilang karena melakukan dua hal buruk sekaligus tanpa terlihat: kendala tenaga
kerja tetap hilang \emph{dan} fungsi tujuannya kini salah, tetapi setiap
pemeriksaan status tetap lolos. Baris yang hilang setidaknya tidak mengubah
fungsi tujuan. Kebiasaan defensif: cetak \texttt{prob} (atau
\texttt{prob.objective}) sebelum menyelesaikan model, dan perlakukan setiap
\texttt{UserWarning} dari PuLP sebagai galat.
\end{enumerate}

\begin{qaflag}
Perbaikan buku diterapkan (catatan sinkronisasi): kotak peringatan dalam
contoh bauran produk (Subbagian 12.12.3, Langkah 4) menyatakan bahwa PuLP
menimpa fungsi tujuan \emph{tanpa pemberitahuan}. Pengujian dengan PuLP 3.3.2
memastikan bahwa \texttt{UserWarning} dicetak. Teks buku diperbarui untuk
menyatakan bahwa peringatan dicetak, tetapi model tetap dapat diselesaikan.
Redaksi latihan itu sendiri tidak perlu diubah; bagian (1) latihan ini merujuk
pada kasus baris yang hilang, yang memang benar-benar tanpa pemberitahuan.
\end{qaflag}

% ----------------------------------------------------------------------------
\exsol{12.8}{Meminimumkan atau Memaksimumkan?}{2}

\begin{enumerate}
\item Meminimumkan $3x_1 + 2x_2$ pada daerah layak menghasilkan
$(x_1, x_2) = (0, 0)$ dengan nilai tujuan $0$ dan status
\texttt{Optimal}. Terverifikasi:

\begin{verbatim}
from pulp import *
# Arah optimisasi salah
prob = LpProblem("Product_Mix", LpMinimize)
x1 = LpVariable("X1", lowBound=0)
x2 = LpVariable("X2", lowBound=0)
prob += 3*x1 + 2*x2, "Total_Profit"
prob += 10*x1 + 5*x2 <= 300, "Material"
prob += 4*x1 + 4*x2 <= 160, "Labor"
prob += 2*x1 + 6*x2 <= 180, "Machine"
prob.solve(PULP_CBC_CMD(msg=0))
print(
    LpStatus[prob.status], value(prob.objective),
    x1.varValue, x2.varValue,
)
\end{verbatim}

Keluaran: \texttt{Optimal 0.0 0.0 0.0}. Model ini tidak taklayak karena semua
kendalanya berupa kendala $\leq$ dengan data nonnegatif: titik asal memenuhi
semuanya. Meminimumkan fungsi tujuan nonnegatif pada daerah yang memuat titik
asal memberikan jawaban yang tidak berguna tetapi sepenuhnya sah, yaitu
``jangan memproduksi apa pun.''

\item Dua perbaikan satu baris, keduanya terverifikasi menghasilkan rencana
$(20, 20)$:
\begin{itemize}
\item Ubah arah optimisasi:
\texttt{prob = LpProblem("Product\_Mix", LpMaximize)}. \\
Keluaran: \texttt{Optimal 100.0 20.0 20.0}.
\item Ubah fungsi tujuan (pertahankan \texttt{LpMinimize}): minimumkan
negatif laba, \texttt{prob += -3*x1 - 2*x2}. \\
Keluaran: \texttt{Optimal -100.0 20.0 20.0}; rencana optimalnya tetap sama,
$(20,20)$, dan laba adalah negatif dari nilai tujuan yang dilaporkan, yaitu
$\$100$.
\end{itemize}

\item \texttt{LpStatus[prob.status]} hanya melaporkan apakah pemecah berhasil
menyelesaikan \emph{model yang Anda bangun}; model dengan arah yang salah
merupakan LP sah dengan optimum yang sah pula, sehingga statusnya
\texttt{Optimal}, baik pada eksekusi yang berkutu maupun yang benar.
Pemeriksaan kewajaran yang berguna: tegaskan bahwa nilai tujuan benar-benar
positif (laba \$0 dari perusahaan yang ingin menghasilkan uang patut
dicurigai); tegaskan bahwa solusi tidak identik nol; bandingkan fungsi tujuan
dengan rencana layak yang diketahui (misalnya, $(10,10)$ menghasilkan \$50,
sehingga optimum setidaknya harus \$50); atau periksa bahwa setidaknya satu
kendala sumber daya mengikat, sebab maksimum laba ketika setiap sumber daya
memiliki slack merupakan hal yang mencurigakan.
\end{enumerate}

% ----------------------------------------------------------------------------
\exsol{12.9}{Laba sebagai Fungsi Tenaga Kerja}{3}

Solusi Terpilih dalam buku sudah benar sebagaimana tercetak; berikut skrip
lengkap dan keluaran sebenarnya (rencana ditambahkan untuk pembahasan bagian 2).

\begin{verbatim}
from pulp import *

def solve_mix(b):
    prob = LpProblem("Bauran_Produk", LpMaximize)
    x1 = LpVariable("X1", lowBound=0)
    x2 = LpVariable("X2", lowBound=0)
    prob += 3*x1 + 2*x2
    prob += 10*x1 + 5*x2 <= 300, "Bahan"
    prob += 4*x1 + 4*x2 <= b, "Tenaga_Kerja"
    prob += 2*x1 + 6*x2 <= 180, "Mesin"
    prob.solve(PULP_CBC_CMD(msg=0))
    laba = value(prob.objective)
    dual = prob.constraints["Tenaga_Kerja"].pi
    return laba, dual, x1.varValue, x2.varValue

for b in range(100, 200, 10):
    laba, dual, v1, v2 = solve_mix(b)
    print(f"b={b}: laba={laba}; dual tenaga kerja={dual}; "
          f"rencana=({v1},{v2})")
\end{verbatim}

Keluaran sebenarnya:

\begin{verbatim}
b=100: laba=75.0; dual tenaga kerja=0.75; rencana=(25.0,0.0)
b=110: laba=82.5; dual tenaga kerja=0.75; rencana=(27.5,0.0)
b=120: laba=90.0; dual tenaga kerja=0.75; rencana=(30.0,0.0)
b=130: laba=92.5; dual tenaga kerja=0.25; rencana=(27.5,5.0)
b=140: laba=95.0; dual tenaga kerja=0.25; rencana=(25.0,10.0)
b=150: laba=97.5; dual tenaga kerja=0.25; rencana=(22.5,15.0)
b=160: laba=100.0; dual tenaga kerja=0.25; rencana=(20.0,20.0)
b=170: laba=102.0; dual tenaga kerja=-0.0; rencana=(18.0,24.0)
b=180: laba=102.0; dual tenaga kerja=-0.0; rencana=(18.0,24.0)
b=190: laba=102.0; dual tenaga kerja=-0.0; rencana=(18.0,24.0)
\end{verbatim}

\begin{enumerate}
\item Kesepuluh labanya adalah
$75, 82.5, 90, 92.5, 95, 97.5, 100, 102, 102, 102$. Kurva tersebut linear
sepotong-sepotong dan cekung, dengan perubahan kemiringan pada $b = 120$ dan
$b = 168$ (nilai terakhir terletak di antara titik tercetak $160$ dan $170$;
di sanalah rencana $(18,24)$, yaitu perpotongan Bahan dan Mesin, menggunakan
tepat $4(18) + 4(24) = 168$ jam tenaga kerja).
\item Kemiringan kurva laba pada setiap $b$ sama dengan harga bayangan
\texttt{c.pi} yang tercetak untuk kendala Tenaga Kerja. Untuk $b \leq 120$,
hanya Tenaga Kerja yang membatasi produksi, rencananya $(b/4, 0)$, dan setiap
jam bernilai $0.75$ (laba naik $7.5$ per $10$ jam:
$75 \to 82.5 \to 90$). Untuk $120 \leq b \leq 168$, Bahan dan Tenaga Kerja
sama-sama mengikat, rencana bergerak sepanjang garis Bahan, dan satu jam
tambahan hanya bernilai $0.25$ (laba naik $2.5$ per $10$ jam). Urutan harga
bayangan yang menurun ketat, $0.75 \to 0.25 \to 0$, tepat mencerminkan
kecekungan kurva.
\item Untuk $b \geq 168$, optimum tertahan pada $(18, 24)$, tempat Bahan dan
Mesin mengikat sementara Tenaga Kerja memiliki slack; tenaga kerja tambahan
memiliki harga bayangan $0$ dan laba tetap datar pada $\$102$. Membeli lebih
banyak sumber daya yang tidak mengikat tidaklah bernilai. Inilah analisis
sensitivitas ruas kanan Bab 10 yang dihitung melalui enumerasi langsung: perulangan
mereproduksi harga bayangan dan rentang yang diizinkannya (dual tetap $0.25$
tepat untuk $120 \leq b \leq 168$) tanpa pengerjaan tableau.
\end{enumerate}
